% GENERATED FILE. Edit canonical lessons, then run npm run generate:books.
% canonical-source-hash: fnv1a64:cff62726e2d76aeb
% canonical-lessons: TE-C361-yuddhabhumi, TE-C361-dalam, TE-C361-muttadincu, TE-C361-jayincu, TE-C361-parajayam

\chapter{On the Battlefield}
\label{ch:te-on-the-battlefield}
\hlchaptermodality{361}

\begin{chapteropening}
\textbf{By the end of this chapter:} \emph{I can say a battlefield, a troop, to besiege, to conquer and a defeat.}

Complete the last lesson of chapter 361: I can say a battlefield, a troop, to besiege, to conquer and a defeat.
\end{chapteropening}

\section[\texorpdfstring{\te{యుద్ధభూమి} (yuddhabhūmi)}{యుద్ధభూమి (yuddhabhūmi)}]{\te{యుద్ధభూమి} (yuddhabhūmi) — a battlefield}

\label{lesson:TE-C361-yuddhabhumi}



\begin{quote}
Before the new one: say the Telugu for a bow, then the Telugu for a gun.
\end{quote}

\subsection*{You'll want to know: \te{యుద్ధభూమి}}
\textbf{\te{యుద్ధభూమి}} — \emph{yuddhabhūmi} — \textquotedblleft{}a battlefield\textquotedblright{}.

\begin{quote}
\te{యుద్ధభూమిలో} \te{రెండు} \te{సైన్యాలు} \te{ఎదురయ్యాయి}.
\end{quote}

\begin{quote}
\emph{yuddhabhūmilō reṇḍu sainyālu edurayyāyi.}
\end{quote}

\begin{quote}
Two armies met on the battlefield.
\end{quote}

\te{యుద్ధం}, war, and \te{భూమి}, ground.

\begin{cousinweb}
\textbf{\te{యుద్ధభూమి}} (\emph{yuddhabhūmi}) is borrowed from Sanskrit: \emph{yuddha}, \textquotedblleft{}battle\textquotedblright{} (CDIAL 10499), and \emph{bhūmi}, \textquotedblleft{}earth, ground\textquotedblright{} (CDIAL 9557).
\end{cousinweb}

\begin{grammarlens}[title={What you've built}]
One more word for telling what happened.
\end{grammarlens}

\subsection*{Your turn}
\begin{itemize}
\raggedright
  \item \emph{Say it:} \emph{yuddhabhūmi}
  \item \emph{Say it:} \emph{yuddhabhūmi}, once more
  \item \emph{Say it:} \emph{yuddhabhūmilō reṇḍu sainyālu edurayyāyi.} — Two armies met on the battlefield.
  \item \emph{Recall:} say the Telugu for a bow, then the Telugu for a gun, then say \emph{yuddhabhūmi} again
\end{itemize}

\begin{tcolorbox}[breakable,colback=teal!4,colframe=teal!35!black,title={Before you move on}]
What does \te{యుద్ధభూమి} mean? (\textquotedblleft{}A battlefield\textquotedblright{}.) Say it once more.
\end{tcolorbox}
\section[\texorpdfstring{\te{దళం} (daḷaṁ)}{దళం (daḷaṁ)}]{\te{దళం} (daḷaṁ) — a troop; a squad}

\label{lesson:TE-C361-dalam}



\begin{quote}
Before the new one: say the Telugu for a gun, then the Telugu for a battlefield.
\end{quote}

\subsection*{You'll want to know: \te{దళం}}
\textbf{\te{దళం}} — \emph{daḷaṁ} — \textquotedblleft{}a troop; a squad\textquotedblright{}.

\begin{quote}
\te{ఒక} \te{చిన్న} \te{దళం} \te{కోటను} \te{కాపాడింది}.
\end{quote}

\begin{quote}
\emph{oka cinna daḷaṁ kōṭanu kāpāḍindi.}
\end{quote}

\begin{quote}
A small troop defended the fort.
\end{quote}

Plural: \te{దళాలు}, \emph{daḷālu}.

Also a petal or a leaf: a part of a whole.

\begin{cousinweb}
\textbf{\te{దళం}} (\emph{daḷaṁ}) is borrowed from Sanskrit \emph{dala}, \textquotedblleft{}a party, a band\textquotedblright{} (CDIAL 6215); CDIAL notes that this sense may come from Dravidian.
\end{cousinweb}

\begin{grammarlens}[title={What you've built}]
One more word for telling what happened.
\end{grammarlens}

\subsection*{Your turn}
\begin{itemize}
\raggedright
  \item \emph{Say it:} \emph{daḷaṁ}
  \item \emph{Say it:} \emph{daḷaṁ}, once more
  \item \emph{Say it:} \emph{oka cinna daḷaṁ kōṭanu kāpāḍindi.} — A small troop defended the fort.
  \item \emph{Recall:} say the Telugu for a gun, then the Telugu for a battlefield, then say \emph{daḷaṁ} again
\end{itemize}

\begin{tcolorbox}[breakable,colback=teal!4,colframe=teal!35!black,title={Before you move on}]
What does \te{దళం} mean? (\textquotedblleft{}A troop; a squad\textquotedblright{}.) Say it once more.
\end{tcolorbox}
\section[\texorpdfstring{\te{ముట్టడించు} (muṭṭaḍiñcu)}{ముట్టడించు (muṭṭaḍiñcu)}]{\te{ముట్టడించు} (muṭṭaḍiñcu) — to besiege; to lay siege}

\label{lesson:TE-C361-muttadincu}



\begin{quote}
Before the new one: say the Telugu for a battlefield, then the Telugu for a troop.
\end{quote}

\subsection*{You'll want to know: \te{ముట్టడించు}}
\textbf{\te{ముట్టడించు}} — \emph{muṭṭaḍiñcu} — \textquotedblleft{}to besiege; to lay siege\textquotedblright{}.

\begin{quote}
\te{శత్రువులు} \te{కోటను} \te{నెల} \te{రోజులు} \te{ముట్టడించారు}.
\end{quote}

\begin{quote}
\emph{śatruvulu kōṭanu nela rōjulu muṭṭaḍiñcāru.}
\end{quote}

\begin{quote}
The enemies besieged the fort for a month.
\end{quote}

Past (I …): \te{ముట్టడించాను}, \emph{muṭṭaḍiñcānu}.

The noun is \te{ముట్టడి}, a siege.

\begin{cousinweb}
\textbf{\te{ముట్టడించు}} (\emph{muṭṭaḍiñcu}) is native Dravidian: DEDR 5018 lists Telugu \emph{muṭṭaḍincu}, \textquotedblleft{}to besiege\textquotedblright{}, and \emph{muṭṭaḍi}, \textquotedblleft{}a siege\textquotedblright{}, with Tamil \emph{muṟṟu}, \textquotedblleft{}to surround, besiege\textquotedblright{}.
\end{cousinweb}

\begin{grammarlens}[title={What you've built}]
One more word for telling what happened.
\end{grammarlens}

\subsection*{Your turn}
\begin{itemize}
\raggedright
  \item \emph{Say it:} \emph{muṭṭaḍiñcu}
  \item \emph{Say it:} \emph{muṭṭaḍiñcu}, once more
  \item \emph{Say it:} \emph{śatruvulu kōṭanu nela rōjulu muṭṭaḍiñcāru.} — The enemies besieged the fort for a month.
  \item \emph{Recall:} say the Telugu for a battlefield, then the Telugu for a troop, then say \emph{muṭṭaḍiñcu} again
\end{itemize}

\begin{tcolorbox}[breakable,colback=teal!4,colframe=teal!35!black,title={Before you move on}]
What does \te{ముట్టడించు} mean? (\textquotedblleft{}To besiege; to lay siege\textquotedblright{}.) Say it once more.
\end{tcolorbox}
\section[\texorpdfstring{\te{జయించు} (jayiñcu)}{జయించు (jayiñcu)}]{\te{జయించు} (jayiñcu) — to conquer}

\label{lesson:TE-C361-jayincu}



\begin{quote}
Before the new one: say the Telugu for a troop, then the Telugu for to besiege.
\end{quote}

\subsection*{You'll want to know: \te{జయించు}}
\textbf{\te{జయించు}} — \emph{jayiñcu} — \textquotedblleft{}to conquer\textquotedblright{}.

\begin{quote}
\te{ఆ} \te{రాజు} \te{చాలా} \te{రాజ్యాలను} \te{జయించాడు}.
\end{quote}

\begin{quote}
\emph{ā rāju cālā rājyālanu jayiñcāḍu.}
\end{quote}

\begin{quote}
That king conquered many kingdoms.
\end{quote}

Past (I …): \te{జయించాను}, \emph{jayiñcānu}.

Also to overcome: \te{భయాన్ని} \te{జయించు}, to conquer fear.

\begin{cousinweb}
\textbf{\te{జయించు}} (\emph{jayiñcu}) is built on Sanskrit \emph{jaya}, \textquotedblleft{}victory\textquotedblright{} (CDIAL 5141), with the Telugu verb ending \emph{-iñcu}.
\end{cousinweb}

\begin{grammarlens}[title={What you've built}]
One more word for telling what happened.
\end{grammarlens}

\subsection*{Your turn}
\begin{itemize}
\raggedright
  \item \emph{Say it:} \emph{jayiñcu}
  \item \emph{Say it:} \emph{jayiñcu}, once more
  \item \emph{Say it:} \emph{ā rāju cālā rājyālanu jayiñcāḍu.} — That king conquered many kingdoms.
  \item \emph{Recall:} say the Telugu for a troop, then the Telugu for to besiege, then say \emph{jayiñcu} again
\end{itemize}

\begin{tcolorbox}[breakable,colback=teal!4,colframe=teal!35!black,title={Before you move on}]
What does \te{జయించు} mean? (\textquotedblleft{}To conquer\textquotedblright{}.) Say it once more.
\end{tcolorbox}
\section[\texorpdfstring{\te{పరాజయం} (parājayaṁ)}{పరాజయం (parājayaṁ)}]{\te{పరాజయం} (parājayaṁ) — a defeat}

\label{lesson:TE-C361-parajayam}



\begin{quote}
Before the new one: say the Telugu for to besiege, then the Telugu for to conquer.
\end{quote}

\subsection*{You'll want to know: \te{పరాజయం}}
\textbf{\te{పరాజయం}} — \emph{parājayaṁ} — \textquotedblleft{}a defeat\textquotedblright{}.

\begin{quote}
\te{ఆ} \te{యుద్ధంలో} \te{రాజుకు} \te{పరాజయం} \te{ఎదురైంది}.
\end{quote}

\begin{quote}
\emph{ā yuddhaṁlō rājuku parājayaṁ eduraindi.}
\end{quote}

\begin{quote}
The king met with defeat in that war.
\end{quote}

The opposite of \te{విజయం}, a victory.

\begin{cousinweb}
\textbf{\te{పరాజయం}} (\emph{parājayaṁ}) is borrowed from Sanskrit \emph{parājaya}, \textquotedblleft{}conquest, defeat\textquotedblright{} (CDIAL 7799p).
\end{cousinweb}

\begin{grammarlens}[title={What you've built}]
One more word for telling what happened.
\end{grammarlens}

\subsection*{Your turn}
\begin{itemize}
\raggedright
  \item \emph{Say it:} \emph{parājayaṁ}
  \item \emph{Say it:} \emph{parājayaṁ}, once more
  \item \emph{Say it:} \emph{ā yuddhaṁlō rājuku parājayaṁ eduraindi.} — The king met with defeat in that war.
  \item \emph{Recall:} say the Telugu for to besiege, then the Telugu for to conquer, then say \emph{parājayaṁ} again
\end{itemize}

\begin{tcolorbox}[breakable,colback=teal!4,colframe=teal!35!black,title={Before you move on}]
What does \te{పరాజయం} mean? (\textquotedblleft{}A defeat\textquotedblright{}.) Say it once more.
\end{tcolorbox}
